\documentclass[11pt]{article}
\usepackage[margin=1in]{geometry}
\usepackage{amsmath,amssymb}
\usepackage{graphicx}
\usepackage{booktabs}
\usepackage[hidelinks]{hyperref}
\usepackage{xcolor}

\title{\bf The MOND Acceleration Scale as the Cosmological Constant:\\
a Coefficient-Free Bridge Between Galaxy Dynamics and Dark-Energy Evolution}
\author{C.\ Zimmerman\\ \small Briar Creek Tech\\ \small\texttt{carlpzimmerman@gmail.com}}
\date{June 2026}

\newcommand{\az}{a_0}
\newcommand{\rde}{\rho_{\rm DE}}

\begin{document}
\maketitle

\begin{abstract}
\noindent
The MOND acceleration scale $\az$ and the cosmological constant $\Lambda$ are numerically related by
$\az \approx c^2\sqrt{\Lambda/32\pi}$, a coincidence noted in various forms by Milgrom, Blanchet, and others. We
take this relation literally---$\az$ \emph{is} the cosmological constant expressed as a galactic acceleration---and
draw out its sharpest consequence. Read as a relation to the dark energy (not the total density), the relation
predicts that the galactically-measured $\az(z)$ and the cosmologically-measured $\sqrt{\rde(z)}$ are the
\emph{same curve}: $\az(z)/\az(0)=\sqrt{\rde(z)/\rde(0)}$, a one-to-one bridge between two independent measurements
in which the (route-forced, undetermined) $32\pi$ coefficient \emph{cancels}. This decouples the framework's
empirical test from its hardest theoretical debt. We exhibit a concrete Aether-Scalar-Tensor (AeST) realization that
carries $\az\propto\sqrt{\rde}$ and verify it is ghost-free, has $c_{\rm GW}=c$, is CMB-safe at linear order, and
keeps $\az$ universal across galaxies. We specify the decisive measurement---a $\sim$30--80-galaxy JWST$+$ALMA
campaign at $z\sim3$---and state the framework's honest limits: the coefficient is traceable but not uniquely
derived, the relation reproduces the value of $\az$ only to $\sim$30\%, the covariant coupling is inserted rather
than forced, the distinctive ($\az$-declining) branch is contingent on DESI's dynamical-dark-energy preference, and
the framework is a unification of the dark sector, \emph{not} a theory of everything. Under a constant $\Lambda$ the
prediction is constant $\az$ (well-supported MOND with its value explained); under DESI's evolving dark energy it is
a falsifiable, mild decline ($\az\!\approx\!0.74\,\az(0)$ at $z=3$) testable now in principle and decisively at
$z\sim3$.
\end{abstract}

\section{Introduction}
The dark sector is $\sim$95\% of the cosmic energy budget: $\sim$68\% dark energy and $\sim$27\% ``dark matter.''
Modified Newtonian Dynamics \cite{Milgrom1983,BekensteinMilgrom1984} replaces the galactic dark-matter
\emph{phenomenology} with a modification of gravity below a universal acceleration $\az\approx1.2\times10^{-10}\,
{\rm m\,s^{-2}}$, reproducing rotation curves and the radial-acceleration relation (RAR) \cite{McGaugh2016} from
the baryons alone with a single constant. It has long been observed \cite{Milgrom1999,Milgrom2009} that this
constant is numerically tied to the cosmological constant,
\begin{equation}
\az \;\approx\; c^2\sqrt{\frac{\Lambda}{32\pi}}\;=\;\frac{c}{2}\sqrt{G\,\rde}\,,
\label{eq:a0Lambda}
\end{equation}
with $\rde=\Lambda c^2/8\pi G$ the dark-energy density. Equation~\eqref{eq:a0Lambda} reproduces the observed $\az$
to $\sim$30\%. The relation has genuine precedents: Milgrom's ``MOND as a vacuum effect'' \cite{Milgrom1999} and
$\az\sim c\sqrt{\Lambda/3}$ coincidence; Blanchet's dipolar dark fluid \cite{Blanchet2009}, in which $\Lambda$ is a
potential floor and $\az$ an anharmonic coefficient; the Cohen--Kaplan--Nelson UV/IR bound \cite{CKN1999}, which
yields $\rde\sim M_P^2 H^2$ of the observed order; and, most recently and most completely, Singh's relativistic
MOND \cite{Singh2026}, in which the single scale is fixed by a dynamically selected de Sitter radius. The deep-MOND
limit itself possesses an exact $SO(4,1)$ (de Sitter) conformal symmetry in three spatial dimensions
\cite{Milgrom2009}; $\az$ is the scale at which that symmetry is \emph{explicitly} broken.

This note does not claim to derive \eqref{eq:a0Lambda}; the $32\pi$ coefficient is traceable to General Relativity
(the Einstein coupling $8\pi$ and the Friedmann $3$) but is route-dependent at the level of a convention, and no
mechanism forces it (Sec.~\ref{sec:honest}). Our contribution is to read \eqref{eq:a0Lambda} as a statement about
the \emph{dark energy}, derive the coefficient-free bridge it implies (Sec.~\ref{sec:bridge}), exhibit a concrete
covariant theory that carries it through the CMB (Sec.~\ref{sec:aest}), and specify the measurement that decides it
(Sec.~\ref{sec:test}).

\section{The two dark components as one scale}
Taken literally, \eqref{eq:a0Lambda} unifies the dark sector: the $\sim$68\% dark energy \emph{is} $\Lambda$, and
the galactic dark-matter phenomenology is the $\az$-modification of gravity sourced by the \emph{same} $\Lambda$.
Dark matter and dark energy become two faces of the vacuum---95\% of the energy budget from a single scale. This is
a unification of the dark sector only: the Standard Model (the remaining $\sim$5\%) is not addressed and is not
derivable from the relation. We return to this honestly in Sec.~\ref{sec:honest}.

\section{The coefficient-free bridge}
\label{sec:bridge}
A subtlety resolves a long-standing ambiguity. The relation has two forms that agree today but diverge in the past:
\begin{equation}
\underbrace{\az=c^2\sqrt{\Lambda/32\pi}=\tfrac{c}{2}\sqrt{G\rde}}_{\propto\,\sqrt{\rde}}
\qquad\text{vs.}\qquad
\underbrace{\az=cH(z)/Z,\ \ H^2=\tfrac{8\pi G}{3}\rho_{\rm tot}}_{\propto\,\sqrt{\rho_{\rm tot}}},
\quad Z=2\sqrt{8\pi/3}.
\end{equation}
They coincide at $z=0$ via the ``why-now'' coincidence $cH_0\approx c^2\sqrt\Lambda$, but $\rho_{\rm tot}$ is
matter-dominated in the past (it rises as $(1+z)^3$) while $\rde$ is not. The form faithful to ``$\az$ is set by the
cosmological constant'' is $\az\propto\sqrt{\rde}$; the $\sqrt{\rho_{\rm tot}}$ form is the apparent-horizon
packaging, and it is disfavoured by existing Tully--Fisher data \cite{Limbach2008,Milgrom2017} (which test exactly
these two couplings and marginally prefer $\sqrt{\rde}$).

Adopting $\az\propto\sqrt{\rde}$, the ratio is independent of the undetermined coefficient:
\begin{equation}
\boxed{\;\frac{\az(z)}{\az(0)} \;=\; \sqrt{\frac{\rde(z)}{\rde(0)}}\;}
\label{eq:bridge}
\end{equation}
since $(c/2)\sqrt{G}$ and $32\pi$ cancel. Equation~\eqref{eq:bridge} is a one-to-one link between the
\emph{galactically}-measured evolution of $\az$ and the \emph{cosmologically}-measured dark-energy density. It is
falsifiable without ever determining $32\pi$:

\begin{center}
\begin{tabular}{lll}
\toprule
cosmology says & galaxies say & verdict\\
\midrule
$\Lambda$ constant ($w=-1$) & $\az$ constant & consistent (the safe core)\\
$\Lambda$ constant ($w=-1$) & $\az$ evolves & \textbf{falsified}\\
dynamical DE (DESI) & $\az\propto\sqrt{\rde}$ & \textbf{confirmed} (distinctive)\\
dynamical DE (DESI) & $\az$ constant & \textbf{falsified}\\
\bottomrule
\end{tabular}
\end{center}

With the DESI DR2 \cite{DESI2025} CPL parameters $(w_0,w_a)=(-0.752,-0.86)$, \eqref{eq:bridge} predicts
$\az(z)/\az(0)=\{1.01,\,0.86,\,0.74\}$ at $z=\{1,2,3\}$---a mild \emph{decline}, the only reading below unity, and
consistent with the high-$z$ Tully--Fisher data that exclude the steep $cH(z)$ rise.

\section{A covariant realization}
\label{sec:aest}
Aether-Scalar-Tensor gravity (AeST) \cite{SkordisZlosnik2021}---a metric, a unit-timelike vector $A^\mu$, and a
scalar $\phi$, governed by a free function $\mathcal{F}(\mathcal{Q},\mathcal{Y})$ of the cosmological invariant
$\mathcal{Q}\sim A^\mu\partial_\mu\phi$ and the spatial-gradient invariant $\mathcal{Y}$---reproduces MOND and fits
the CMB. We realize $\az\propto\sqrt{\rde}$ with the non-separable free function
\begin{equation}
\mathcal{F}(\mathcal{Q},\mathcal{Y}) = K(\mathcal{Q}) + C(\mathcal{Q})\,\mathcal{Y}^{3/2} + (\text{Newtonian
crossover}),\qquad C(\mathcal{Q}) = \kappa_0\sqrt{\rho_0/\rde(\mathcal{Q})}.
\end{equation}
Because the deep-MOND scale obeys $\az\propto1/(\text{the }\mathcal{Y}^{3/2}\text{ coefficient})$, this gives
$\az(\mathcal{Q})=\az(0)\sqrt{\rde(\mathcal{Q})/\rho_0}$ exactly: the bridge \eqref{eq:bridge} falls out of the
action. Four checks pass (verified symbolically; see the accompanying repository):
\begin{itemize}\itemsep1pt
\item \textbf{Ghost-free / hyperbolic:} $\mathcal{F}_{\mathcal{Y}}=\tfrac32 C\sqrt{\mathcal{Y}}$ and
$2\mathcal{Y}\mathcal{F}_{\mathcal{Y}\mathcal{Y}}+\mathcal{F}_{\mathcal{Y}}=3C\sqrt{\mathcal{Y}}$ are both $>0$
iff $C>0$, which holds wherever $\rde>0$.
\item \textbf{$c_{\rm GW}=c$:} $\mathcal{F}$ is a scalar potential; it does not enter the graviton kinetic term
(consistent with GW170817).
\item \textbf{CMB-safe:} $\mathcal{Y}$ is gradient-built, so $\bar{\mathcal{Y}}=0$ on FRW and
$C(\mathcal{Q})\mathcal{Y}^{3/2}=\mathcal{O}(\delta^3)$; $\az$ is absent from the linear equations regardless of its
value, and tying $C$ to the \emph{background} $\mathcal{Q}$ does not promote it.
\item \textbf{Universality:} $\mathcal{Q}$ is the cosmological time-sector invariant; galaxies source the spatial
$\mathcal{Y}$ (the MOND potential) and perturb $\mathcal{Q}$ only at the potential level $\Phi=(V/c)^2\sim10^{-6}$,
so $\delta\az/\az\sim10^{-6}$--$10^{-3}\ll$ the RAR scatter ($\sim$29\%). (Order-of-magnitude; a full quasi-static
solution remains.)
\end{itemize}
The $C(\mathcal{Q})\to\infty$ strong-coupling limit (dynamical DE, $\rde\to0$ in the deep past) is benign: it occurs
at $z\sim1100$ where no nonlinear structures exist ($\mathcal{Y}\!\approx\!0$), while the first galaxies form at
$z\lesssim10$--20 where $C$ is only $\sim$3--5$\times$; the two epochs do not overlap, $\az\to0$ smoothly, and it is a
continuation of deep-MOND's intrinsic $\mathcal{Y}^{3/2}$ strong coupling shared by all MOND theories.

\section{The decisive measurement}
\label{sec:test}
The bridge \eqref{eq:bridge} is tested via the baryonic Tully--Fisher zero point: $M_b=V_{\rm flat}^4/(G\az)$, so at
fixed $M_b$, $V_{\rm flat}\propto\az^{1/4}$. The coefficient-free ratio $\az(z)/\az(0)$ is read from the shift of
the BTFR zero point between a $z\sim3$ sample and a $z\sim0$ anchor analyzed through an identical pipeline (so the
IMF, stellar-mass, interpolation-function, and absolute-$\az$ systematics cancel). The faithful-vs-constant signal
at $z=3$ is a $0.033$\,dex shift in $V_{\rm flat}$; the excluded $cH(z)$ rise ($+0.16$\,dex) is trivially separated.
A forecast with per-galaxy precision $\sim$0.06\,dex gives a $3\sigma$ ($5\sigma$) separation with
$\sim$30 ($\sim$80) low-surface-density, rotation-supported deep-MOND discs at $z=2.5$--3.5, using JWST/NIRSpec
outer rotation curves, ALMA [C\,{\sc ii}]/CO gas masses, and JWST SED stellar masses. The measurement requires no
new facility and no determination of the coefficient. (DESI itself supplies the $\rde(z)$ input but, being a
redshift survey, cannot measure $\az$'s evolution; its peculiar-velocity component is $z<0.15$.)

\section{Honest assessment}
\label{sec:honest}
We state the limits plainly, since they bound what \eqref{eq:a0Lambda}--\eqref{eq:bridge} can claim.
\begin{enumerate}\itemsep1pt
\item \textbf{Not novel at the core.} The $\az$--$\Lambda$ relation and the deep-MOND $SO(4,1)$/de Sitter structure
are Milgrom's \cite{Milgrom1999,Milgrom2009}; the most complete relativistic synthesis is Singh's \cite{Singh2026}.
The contributions here are the coefficient-free \emph{bridge} framing \eqref{eq:bridge} and the explicit AeST
realization with its four verified checks.
\item \textbf{The coefficient is not derived.} $32\pi=4\times8\pi$ is GR-traceable (the Einstein coupling, the
Friedmann $3$) but route-forced; no symmetry, bound, or dynamical mechanism we examined uniquely fixes it. It enters
only the $\sim$30\% \emph{value} check, not the bridge.
\item \textbf{The covariant coupling is inserted.} $C(\mathcal{Q})$ is consistent and healthy but put in by hand; no
symmetry selects it.
\item \textbf{Contingent and modest.} The distinctive declining branch is contingent on DESI's dynamical-DE
preference ($\sim$3--4$\sigma$, contested). Under a true constant $\Lambda$ the prediction is simply constant $\az$.
\item \textbf{A theory of the dark universe, not of everything.} The Standard Model is input; gravity is taken to be
emergent; the value of $\Lambda$ (the cosmological-constant problem) and a UV completion of MOND's strong coupling
remain open---as they do for all competitors.
\end{enumerate}

\section{Conclusion}
If the MOND scale is the cosmological constant, then galaxy dynamics and dark-energy evolution are two views of one
quantity, related by the coefficient-free bridge $\az(z)=\az(0)\sqrt{\rde(z)/\rde(0)}$. The relation is realized in a
concrete, ghost-free, CMB- and GW-safe covariant theory; its distinctive content is decided by a single
$\sim$30--80-galaxy measurement at $z\sim3$; and its honest status is a unification of the dark sector that survives
present consistency tests while leaving its coefficient, its UV completion, and the Standard Model untouched. We
advocate stating it as exactly that.

\section*{Data and code availability}
The verification scripts (symbolic ghost/CMB/$c_{\rm GW}$/universality checks, the bridge and forecast
calculations) are available at \url{https://github.com/carlzimmerman/zimmerman-formula}.

\section*{Acknowledgement of methods}
Symbolic and numerical verification, literature cross-checking, and drafting were performed with AI assistance
(Anthropic Claude); all physical claims and their honest limitations were checked against primary sources and are
the author's responsibility.

\begin{thebibliography}{99}
\small
\bibitem{Milgrom1983} M.\ Milgrom, \emph{ApJ} \textbf{270}, 365 (1983).
\bibitem{BekensteinMilgrom1984} J.\ Bekenstein and M.\ Milgrom, \emph{ApJ} \textbf{286}, 7 (1984).
\bibitem{McGaugh2016} S.\ McGaugh, F.\ Lelli, J.\ Schombert, \emph{PRL} \textbf{117}, 201101 (2016).
\bibitem{Milgrom1999} M.\ Milgrom, \emph{Phys.\ Lett.\ A} \textbf{253}, 273 (1999), astro-ph/9805346.
\bibitem{Milgrom2009} M.\ Milgrom, \emph{ApJ} \textbf{698}, 1630 (2009), arXiv:0810.4065.
\bibitem{Blanchet2009} L.\ Blanchet and A.\ Le Tiec, \emph{Phys.\ Rev.\ D} \textbf{80}, 023524 (2009),
arXiv:0901.3114.
\bibitem{CKN1999} A.\ Cohen, D.\ Kaplan, A.\ Nelson, \emph{PRL} \textbf{82}, 4971 (1999), hep-th/9803132.
\bibitem{Li2004} M.\ Li, \emph{Phys.\ Lett.\ B} \textbf{603}, 1 (2004), hep-th/0403127.
\bibitem{Singh2026} T.\ P.\ Singh, ``A Relativistic MOND,'' arXiv:2601.04290 (2026).
\bibitem{SkordisZlosnik2021} C.\ Skordis and T.\ Z\l{}o\'snik, \emph{PRL} \textbf{127}, 161302 (2021),
arXiv:2007.00082.
\bibitem{Limbach2008} S.\ Limbach, D.\ Psaltis, F.\ \"Ozel, arXiv:0809.2790 (2008).
\bibitem{Milgrom2017} M.\ Milgrom, ``High-redshift rotation curves and MOND,'' arXiv:1703.06110 (2017).
\bibitem{DESI2024} DESI Collaboration, \emph{Phys.\ Rev.\ D} \textbf{111}, 023532 (2025), arXiv:2404.03002.
\bibitem{DESI2025} DESI Collaboration, ``DESI DR2 Results II,'' arXiv:2503.14738 (2025).
\end{thebibliography}

\end{document}
